% GENERATED FILE. Edit canonical lessons, then run npm run generate:books.
% canonical-source-hash: fnv1a64:6720a1b9a856f6c1
% canonical-lessons: ZH-C237-zhiwu, ZH-C237-huasheng, ZH-C237-yezi, ZH-C237-faya, ZH-C237-xigua

\chapter{Plants}
\label{ch:zh-plants}
\hlchaptermodality{237}

\begin{chapteropening}
\textbf{By the end of this chapter:} \emph{I can talk about plants, peanuts, leaves, things sprouting and watermelons.}

Complete the last lesson of chapter 237: I can talk about plants, peanuts, leaves, things sprouting and watermelons.
\end{chapteropening}

\section[zhíwù]{zhíwù — plant}

\label{lesson:ZH-C237-zhiwu}



\begin{quote}
Before the new one: say the Mandarin for luckily, then the Mandarin for nearly.
\end{quote}

\subsection*{You'll want to know: zhíwù}
\emph{zhíwù} — \textquotedblleft{}plant\textquotedblright{}.

\begin{quote}
\emph{wǒ xǐhuan zhíwù} — I love plants.
\end{quote}

Its characters come later; for now it is a word for the ear.

\begin{cousinweb}
\emph{zhíwù} is built from \emph{zhí} 'to plant' and \emph{wù} 'thing', both inherited from Middle Chinese, \emph{dzyik} and \emph{mjut} (Guangyun, in Baxter's transcription), which glosses the first 'to plant' and 'to set up'.
\end{cousinweb}

\begin{grammarlens}[title={What you've built}]
One more word for everyday things.
\end{grammarlens}

\subsection*{Your turn}
\begin{itemize}
\raggedright
  \item \emph{Say it:} \emph{zhíwù}
  \item \emph{Say it:} \emph{zhíwù}, once more
  \item \emph{Say it:} \emph{wǒ xǐhuan zhíwù} — I love plants.
  \item \emph{Recall:} say the Mandarin for luckily, then the Mandarin for nearly, then say \emph{zhíwù} again
\end{itemize}

\begin{tcolorbox}[breakable,colback=teal!4,colframe=teal!35!black,title={Before you move on}]
What does zhíwù mean? (\textquotedblleft{}Plant\textquotedblright{}.) Say it once more.
\end{tcolorbox}
\section[\texorpdfstring{\zh{花生} (huāshēng)}{花生 (huāshēng)}]{\zh{花生} (huāshēng) — peanut}

\label{lesson:ZH-C237-huasheng}



\begin{quote}
Before the new one: say the Mandarin for nearly, then the Mandarin for a plant.
\end{quote}

\subsection*{You'll want to know: \zh{花生}}
\textbf{\zh{花生}} — \emph{huāshēng} — \textquotedblleft{}peanut\textquotedblright{}.

\begin{quote}
\textbf{\zh{我现在不吃花生了}} \emph{wǒ xiànzài bù chī huāshēng le} — I don't eat peanuts any more.
\end{quote}

\begin{cousinweb}
\textbf{\zh{花生}} (\emph{huāshēng}) is built from \zh{花} 'flower' and \zh{生} 'to grow', both inherited from Middle Chinese, \emph{xwae} and \emph{srjaeng} (Guangyun, in Baxter's transcription); how the pair came to name the peanut is not settled.
\end{cousinweb}

\begin{grammarlens}[title={What you've built}]
One more word for everyday things.
\end{grammarlens}

\subsection*{Your turn}
\begin{itemize}
\raggedright
  \item \emph{Say it:} \emph{huāshēng}
  \item \emph{Say it:} \emph{huāshēng}, once more
  \item \emph{Say it:} \emph{wǒ xiànzài bù chī huāshēng le} — I don't eat peanuts any more.
  \item \emph{Recall:} say the Mandarin for nearly, then the Mandarin for a plant, then say \emph{huāshēng} again
\end{itemize}

\begin{tcolorbox}[breakable,colback=teal!4,colframe=teal!35!black,title={Before you move on}]
What does \zh{花生} mean? (\textquotedblleft{}Peanut\textquotedblright{}.) Say it once more.
\end{tcolorbox}
\section[yèzi]{yèzi — leaf}

\label{lesson:ZH-C237-yezi}



\begin{quote}
Before the new one: say the Mandarin for a plant, then the Mandarin for a peanut.
\end{quote}

\subsection*{You'll want to know: yèzi}
\emph{yèzi} — \textquotedblleft{}leaf\textquotedblright{}.

\begin{quote}
\emph{zhèxiē yèzi hěn piàoliang} — These leaves are lovely.
\end{quote}

Its characters come later; for now it is a word for the ear.

\begin{cousinweb}
\emph{yèzi} is built from \emph{yè} 'leaf' and the noun ending \zh{子}, both inherited from Middle Chinese, \emph{yep} and \emph{tsiX} (Guangyun, in Baxter's transcription), which glosses the first 'branches and leaves'.
\end{cousinweb}

\begin{grammarlens}[title={What you've built}]
One more word for everyday things.
\end{grammarlens}

\subsection*{Your turn}
\begin{itemize}
\raggedright
  \item \emph{Say it:} \emph{yèzi}
  \item \emph{Say it:} \emph{yèzi}, once more
  \item \emph{Say it:} \emph{zhèxiē yèzi hěn piàoliang} — These leaves are lovely.
  \item \emph{Recall:} say the Mandarin for a plant, then the Mandarin for a peanut, then say \emph{yèzi} again
\end{itemize}

\begin{tcolorbox}[breakable,colback=teal!4,colframe=teal!35!black,title={Before you move on}]
What does yèzi mean? (\textquotedblleft{}Leaf\textquotedblright{}.) Say it once more.
\end{tcolorbox}
\section[fāyá]{fāyá — to sprout; to bud}

\label{lesson:ZH-C237-faya}



\begin{quote}
Before the new one: say the Mandarin for a peanut, then the Mandarin for a leaf.
\end{quote}

\subsection*{You'll want to know: fāyá}
\emph{fāyá} — \textquotedblleft{}to sprout; to bud\textquotedblright{}.

\begin{quote}
\emph{chūntiān lái le xiǎo cǎo fāyá le} — Spring is here, and the grass is coming up.
\end{quote}

Its characters come later; for now it is a word for the ear.

\begin{cousinweb}
\emph{fāyá} is built from \emph{fā} 'to send out' and \emph{yá} 'a shoot'. The first is inherited from Middle Chinese \emph{pjot} (Guangyun, in Baxter's transcription), glossed 'to rise' and 'to open out'; the dictionary glosses the second, \emph{ngae} there, 'a sprout'.
\end{cousinweb}

\begin{grammarlens}[title={What you've built}]
One more word for everyday things.
\end{grammarlens}

\subsection*{Your turn}
\begin{itemize}
\raggedright
  \item \emph{Say it:} \emph{fāyá}
  \item \emph{Say it:} \emph{fāyá}, once more
  \item \emph{Say it:} \emph{chūntiān lái le xiǎo cǎo fāyá le} — Spring is here, and the grass is coming up.
  \item \emph{Recall:} say the Mandarin for a peanut, then the Mandarin for a leaf, then say \emph{fāyá} again
\end{itemize}

\begin{tcolorbox}[breakable,colback=teal!4,colframe=teal!35!black,title={Before you move on}]
What does fāyá mean? (\textquotedblleft{}To sprout; to bud\textquotedblright{}.) Say it once more.
\end{tcolorbox}
\section[xīguā]{xīguā — watermelon}

\label{lesson:ZH-C237-xigua}



\begin{quote}
Before the new one: say the Mandarin for a leaf, then the Mandarin for to sprout.
\end{quote}

\subsection*{You'll want to know: xīguā}
\emph{xīguā} — \textquotedblleft{}watermelon\textquotedblright{}.

\begin{quote}
\emph{xiàtiān wǒmen chángcháng chī xīguā} — We eat a lot of watermelon in summer.
\end{quote}

Its characters come later; for now it is a word for the ear.

\begin{cousinweb}
\emph{xīguā} is built from \zh{西} 'west' and \emph{guā} 'melon', both inherited from Middle Chinese, \emph{sej} and \emph{kwae} (Guangyun, in Baxter's transcription), which glosses the second as a gourd and lists many kinds.
\end{cousinweb}

\begin{grammarlens}[title={What you've built}]
One more word for everyday things.
\end{grammarlens}

\subsection*{Your turn}
\begin{itemize}
\raggedright
  \item \emph{Say it:} \emph{xīguā}
  \item \emph{Say it:} \emph{xīguā}, once more
  \item \emph{Say it:} \emph{xiàtiān wǒmen chángcháng chī xīguā} — We eat a lot of watermelon in summer.
  \item \emph{Recall:} say the Mandarin for a leaf, then the Mandarin for to sprout, then say \emph{xīguā} again
\end{itemize}

\begin{tcolorbox}[breakable,colback=teal!4,colframe=teal!35!black,title={Before you move on}]
What does xīguā mean? (\textquotedblleft{}Watermelon\textquotedblright{}.) Say it once more.
\end{tcolorbox}
